\documentclass[11pt]{article}
\usepackage{amsmath,amssymb,amsthm,mathrsfs}
\usepackage[margin=1in]{geometry}
\usepackage{hyperref}

\newtheorem{theorem}{Theorem}
\newtheorem{proposition}[theorem]{Proposition}
\newtheorem{lemma}[theorem]{Lemma}
\newtheorem{corollary}[theorem]{Corollary}
\theoremstyle{definition}
\newtheorem{remark}[theorem]{Remark}
\newtheorem{convention}[theorem]{Convention}

\DeclareMathOperator{\ord}{ord}
\DeclareMathOperator{\maj}{maj}
\DeclareMathOperator{\comaj}{comaj}
\newcommand{\hgt}{\operatorname{ht}}   % NOT \ht -- collides with a TeX primitive
\newcommand{\Z}{\mathbb{Z}}
\newcommand{\Q}{\mathbb{Q}}
\newcommand{\C}{\mathbb{C}}
\newcommand{\Lam}{\Lambda}
\newcommand{\F}{\mathcal{F}}
\newcommand{\Ad}{\operatorname{Ad}}

\title{The hyperbola $ts=1$ is a fixed-point set --- but not of the involution
we were looking for}
\author{Clio}
\date{9 September 2026}

\begin{document}
\maketitle

\begin{center}\small
\fbox{\parbox{0.92\textwidth}{\centering
\textbf{Author:} Clio \quad\textbullet\quad \textbf{Date:} 9 September 2026 \quad\textbullet\quad
\textbf{Recipient:} Rick (cc Robin)\\
\textbf{Title:} The hyperbola $ts=1$ is a fixed-point set --- but not of the
involution we were looking for\\
\textbf{Corresponds to:} \texttt{clio-vega/proofs@7002430} (parent commit; this paper is the tip)
}}
\end{center}
\bigskip

\begin{abstract}
Let $R_e(t)$ be the operator on $\Lam$ adding a connected $e$-ribbon with weight
$t^{\hgt}$.  Seven sessions have recorded that the locus $ts=1$ in the two-parameter
family $[R_e(t),R_f(s)]$ is distinguished, without saying why.  The proposed
explanation (Q104) was that $ts=1$ is the fixed-point set of a bead-parity involution
$\Theta$ with $\Theta R_e(t)\Theta=R_e(1/t)$.

We decide this.  Three results and one honest negative.

\textbf{(A)} The gating computation Q107 is settled in closed form: if $\lambda/\mu$
is a connected $e$-ribbon $R$ whose topmost row is $r$, then
$n(\lambda)-n(\mu)=(r-1)e+\maj(R)$.  The Garsia--Zabrocki exponent is therefore
\emph{never} constant in $\mu$ --- not for any ribbon shape, columns included ---
and the conjugation is not an operator identity with a scalar.  Q104 \emph{fails as
stated}.  We locate the exact pin that makes the Garsia--Zabrocki theorem work:
their operator adds cells at rows $1,\dots,k$ \emph{fixed independently of $\mu$},
whereas a ribbon floats.  It is the row-pin, not the column shape.

\textbf{(B)} Q104 nevertheless has a true answer, and the involution is the most
classical one available:
\[
  \boxed{\ \omega\,R_e(t)\,\omega \;=\; t^{\,e-1}\,R_e(1/t),\qquad \omega:s_\lambda\mapsto s_{\lambda'}.\ }
\]
The anomaly $t^{e-1}$ is forced by $\hgt(R)+\hgt(R^{\mathsf T})=e-1$.

\textbf{(C)} Consequently $ts=1$ \emph{is} a fixed-point set: of the involution
$\sigma(t,s)=(1/s,1/t)$ on the parameter torus, which is implemented on operators by
$\Ad(\omega)$ composed with the swap of the two commutator slots.  On $ts=1$ the
commutator $[R_e(t),R_e(s)]$ lies in the $(-1)$-eigenspace of $\Ad(\omega)$.

\textbf{(D)} The two-bead sector of $[R_e(t),R_f(s)]$ vanishes on $ts=1$ for
\emph{all} $e,f\ge1$, and the mechanism is local on the Maya diagram: reversing the
order of two disjoint bead moves multiplies the weight by exactly $(ts)^k$.

\textbf{The negative, stated plainly.}  (C) does \emph{not} imply (D).  The one-bead
sector also lies in the $(-1)$-eigenspace on $ts=1$ and is \emph{nonzero} there
($52/52$ matrix entries survive, against $0/58$ in the two-bead sector).  The
programme's hope that all four sightings of $ts=1$ become corollaries of one
$\Z/2$ is therefore \textbf{false}: two of them do, two do not.
\end{abstract}

\section{Conventions and the gate}

\begin{convention}\label{conv:main}
$\Lam$ is the ring of symmetric functions over $\Q(t,s)$ with Schur basis
$\{s_\lambda\}$.  For a partition $\lambda$, $n(\lambda)=\sum_{i\ge1}(i-1)\lambda_i$,
$M(\lambda)=\{\lambda_j-j:j\ge1\}$ is the Maya set, and for $e\ge1$
\[
  R_e(t)\,s_\mu=\sum_{\lambda}t^{\hgt(\lambda/\mu)}s_\lambda,
\]
the sum over $\lambda\supset\mu$ with $\lambda/\mu$ a connected border strip of size
$e$ and $\hgt=(\#\text{rows})-1$.  We write $C_{e,f}(t,s)=[R_e(t),R_f(s)]$ and
$\omega$ for the standard involution $s_\lambda\mapsto s_{\lambda'}$, extended
$\Q(t,s)$-linearly.  This is Convention 2.1 of
\texttt{2026-09-07-Q92-cross-rank-commutator.tex}, unchanged.
\end{convention}

\begin{remark}[the gate, and it failed]\label{rem:gate}
The framing of Q104 carried a flagged dependency: a substitution
$t=-e^{\lambda},\ s=-e^{\mu}$ attributed to Zhang,
\href{https://arxiv.org/abs/2603.11172}{\texttt{2603.11172}}, eq.~(30).  We opened the
paper's \LaTeX{} source and located equation (30) by counting numbered environments.
It is Shastry's decorated Yang--Baxter equation for the spinful Hubbard
$R$-matrix,
\[
  \bigl(\check R^{\uparrow\downarrow}_{23}(\lambda-\mu)
   +f(\lambda,\mu)\check R^{\uparrow\downarrow}_{23}(\lambda+\mu)
   \Theta^\uparrow_2\Theta^\downarrow_2\bigr)
  \check R^{\uparrow\downarrow}_{12}(\lambda)\check R^{\uparrow\downarrow}_{23}(\mu)
  =\cdots,
\]
in which $\lambda,\mu$ are \emph{spectral parameters} and $\Theta_j=2n_j-1$ is a
time-reversal operator on a spin-chain site.  A census over the full source returns
$0$ hits for the exponential substitution and $0$ each for ``Hall--Littlewood'',
``ribbon'', ``Fock'' and ``symmetric function''.  \textbf{The substitution is not in
the paper.}  It was a sub-agent's construction.  Zhang is dropped entirely; nothing
below depends on it, and the involution used below is defined from
Convention~\ref{conv:main}.
\end{remark}

\section{Q107 in closed form}

\begin{convention}[$\maj$ of a border strip]\label{conv:maj}
Let $R$ be a connected border strip of size $e$ occupying $m$ rows with lengths
$a_1,\dots,a_m$ read \emph{top to bottom}.  Read the cells of $R$ starting from its
bottom-left end; each successive cell is one step right or one step up.  The
\emph{descent set} is $D(R)=\{p: \text{the step from cell }p\text{ to cell }p+1
\text{ is up}\}\subseteq\{1,\dots,e-1\}$, so
$D(R)=\{a_m,\ a_m+a_{m-1},\ \dots,\ a_m+\cdots+a_2\}$ and $|D(R)|=m-1=\hgt(R)$.
Put $\maj(R)=\sum_{p\in D(R)}p$.
\end{convention}

\begin{theorem}[Q107]\label{thm:A}
Let $\mu\subset\lambda$ with $\lambda/\mu$ a connected border strip of size $e$ and
shape $R$, whose topmost row is $r$ (rows indexed from $1$).  Then
\[
  \boxed{\ c(R,\mu)\;:=\;n(\lambda)-n(\mu)\;=\;(r-1)\,e\;+\;\maj(R).\ }
\]
\end{theorem}

\begin{proof}
Since $n(\nu)=\sum_{(i,j)\in\nu}(i-1)$, summing over cells,
\[
  n(\lambda)-n(\mu)=\sum_{(i,j)\in\lambda/\mu}(i-1).
\]
The strip occupies rows $r,r+1,\dots,r+m-1$ with $a_1,\dots,a_m$ cells respectively,
$\sum_i a_i=e$.  Hence
\begin{equation}\label{eq:split}
  n(\lambda)-n(\mu)=\sum_{i=1}^{m}a_i\,(r+i-1-1)=(r-1)e+\sum_{i=1}^{m}(i-1)a_i .
\end{equation}
It remains to identify $\sigma(R):=\sum_{i=1}^m(i-1)a_i$ with $\maj(R)$.  Count the
quantity $\Sigma=\sum_{p=1}^{e}\#\{d\in D(R):d<p\}$ in two ways.

Cell-first: a cell in row $i$ (from the top) lies in the $(m-i+1)$-st block of the
bottom-to-top reading, so exactly $m-i$ descents precede it.  Thus
$\Sigma=\sum_{i=1}^m a_i(m-i)$.

Descent-first: each $d\in D(R)$ is counted once for every $p$ with $d<p\le e$, i.e.\
$e-d$ times.  Thus $\Sigma=\sum_{d\in D(R)}(e-d)=\hgt(R)\,e-\maj(R)$, using
$|D(R)|=m-1=\hgt(R)$.

Since $(m-i)+(i-1)=m-1=\hgt(R)$ for every $i$,
\[
  \sum_{i=1}^m a_i(m-i)\;+\;\sum_{i=1}^m a_i(i-1)\;=\;\hgt(R)\sum_i a_i=\hgt(R)\,e,
\]
so $\sigma(R)=\hgt(R)e-\Sigma=\hgt(R)e-\bigl(\hgt(R)e-\maj(R)\bigr)=\maj(R)$.
Substituting into \eqref{eq:split} gives the claim.
\end{proof}

\begin{corollary}[Q107 answers \textup{(b)}, and Q104 fails as stated]\label{cor:notconst}
For every $e\ge1$ and every ribbon shape $R$, $c(R,\mu)$ is \emph{not} constant as
$\mu$ varies: its value set is $\{(r-1)e+\maj(R)\}$ over the attainable top-rows $r$,
an infinite arithmetic progression of common difference $e$.  Its level sets are
exactly the fibres of $r$.  Consequently the Garsia--Zabrocki $t\to1/t$ conjugation
of $R_e(t)$ emits an exponent that varies with the matrix element, and there is no
scalar $c$ with $\Theta R_e(t)\Theta=t^{c}R_e(1/t)$ obtained by that mechanism.
\end{corollary}

\begin{remark}[where the constancy in Garsia--Zabrocki actually comes from]
\label{rem:gz}
Garsia--Zabrocki, \emph{Polynomiality of the $q,t$-Kostka Revisited},
\href{https://arxiv.org/abs/math/0008199}{\texttt{math/0008199}}, Theorem 2.3, is
proved by setting $t\to1/t$ in their (2.11) to obtain
\[
  t^{-n(\mu)}\,\mathcal H_k(X;1/t)\,H_\mu[X;t]
  \;=\;t^{-n(\mu+1^k)}\,H_{\mu+1^k}[X;t],
\]
and then, verbatim, ``\emph{and this (using 2.10) may be rewritten as}''
$t^{\binom k2}\mathcal H_k(X;1/t)H_\mu[X;t]=H_{\mu+1^k}[X;t]$.  The step
``using 2.10'' is exactly the assertion $n(\mu+1^k)-n(\mu)=\binom k2$, i.e.\ our
$c$ for the shape they add.  It is constant --- but the reason is \emph{not} that
$1^k$ is a column.  The operation $\mu\mapsto\mu+1^k$ adds one cell to each of rows
$1,\dots,k$, a set of rows fixed \emph{independently of $\mu$}; by \eqref{eq:split}
the exponent is then $\sum_{i=1}^k(i-1)=\binom k2$ whatever $\mu$ is.  A ribbon has
no such pin: it floats, and $r$ moves.

Two consequences worth recording, both of which correct the reading this session
inherited.
\begin{enumerate}
\item[(i)] The prediction carried in from the browse log of 2026-09-08 --- that $c$
is ``constant in $\mu$ only for columns'' --- is \textbf{false even for columns}.  A
column ribbon $1^k$ added at top-row $r$ has $c=(r-1)k+\binom k2$; for $e=3$ the
shape $(1,1,1)$ realises $c\in\{3,6,9,12,\dots\}$.  The prediction was shape-correct
about \emph{which theorem} was doing the work and wrong about \emph{why}.
\item[(ii)] $\mu\mapsto\mu+1^k$ is in general a \emph{disconnected} skew shape, not a
ribbon, and its $c$ is constant.  So the negative control proposed in the session
brief --- ``break the ribbon condition and check $c$ is not constant there'' ---
tests the wrong variable.  Connectedness is irrelevant; row-pinning is everything.
\end{enumerate}
\end{remark}

\section{Q104: the involution exists, and it is $\omega$}

\begin{lemma}\label{lem:transpose}
Let $\lambda/\mu$ be a connected border strip of size $e$.  Then $\lambda'/\mu'$ is a
connected border strip of size $e$, and
\[
  \hgt(\lambda/\mu)+\hgt(\lambda'/\mu')=e-1 .
\]
\end{lemma}

\begin{proof}
Transposition is an involution on Young diagrams preserving containment, so
$\lambda'/\mu'$ is a skew shape of size $e$; it is connected and contains no $2\times2$
block because both properties are transposition-invariant.  For the height: let the
strip occupy $\rho$ rows and $\gamma$ columns.  Reading its $e$ cells from one end to
the other, the $e-1$ steps between consecutive cells are each either a row-change or a
column-change, and the row-changes number $\rho-1$, the column-changes $\gamma-1$.
Hence $(\rho-1)+(\gamma-1)=e-1$, i.e. $\rho+\gamma=e+1$.  Transposition exchanges rows
and columns, so $\hgt(\lambda'/\mu')=\gamma-1=e-\rho=e-1-(\rho-1)
=e-1-\hgt(\lambda/\mu)$.
\end{proof}

\begin{theorem}[Q104, corrected form]\label{thm:B}
For every $e\ge1$,
\[
  \boxed{\ \omega\,R_e(t)\,\omega\;=\;t^{\,e-1}\,R_e(1/t).\ }
\]
Equivalently, $\widetilde R_e(t):=t^{-(e-1)/2}R_e(t)$ satisfies
$\omega\widetilde R_e(t)\omega=\widetilde R_e(1/t)$.
\end{theorem}

\begin{proof}
Both sides are $\Q(t)$-linear, so it suffices to evaluate on $s_\mu$.  Since
$\omega^2=\mathrm{id}$ and $\omega s_\nu=s_{\nu'}$,
\[
  \omega R_e(t)\omega\, s_\mu
  =\omega R_e(t)\,s_{\mu'}
  =\omega\!\!\sum_{\nu\supset\mu',\ \nu/\mu'\ e\text{-strip}}\!\! t^{\hgt(\nu/\mu')}s_\nu
  =\!\!\sum_{\nu\supset\mu'}\!\! t^{\hgt(\nu/\mu')}s_{\nu'} .
\]
As $\nu$ runs over the partitions with $\nu/\mu'$ a connected $e$-strip,
$\lambda:=\nu'$ runs over the partitions with $\lambda/\mu$ a connected $e$-strip
(Lemma~\ref{lem:transpose} applied in both directions), and
$\hgt(\nu/\mu')=e-1-\hgt(\lambda/\mu)$.  Hence the sum equals
\[
  \sum_{\lambda\supset\mu} t^{\,e-1-\hgt(\lambda/\mu)}s_\lambda
  \;=\;t^{e-1}\!\!\sum_{\lambda\supset\mu}\!(1/t)^{\hgt(\lambda/\mu)}s_\lambda
  \;=\;t^{e-1}R_e(1/t)\,s_\mu . \qedhere
\]
\end{proof}

\begin{remark}[why this was missed for seven sessions]
The browse log of 2026-09-08 recorded, from Rozhkovskaya
\href{https://arxiv.org/abs/2202.01544}{\texttt{2202.01544}} (4.1)--(4.2), that
$\omega\Gamma^+(u)\omega^{-1}=\Gamma^-(-u)$ \emph{fixes $t$}, and concluded
``\,$\Theta$ cannot be $\omega$''.  That conclusion is correct \emph{about
$\Gamma^\pm$} and false about $R_e$.  The two operators carry different $t$-statistics:
Rozhkovskaya's $t$ sits in structure constants $(1-t)$, $(1-t^i)$, whereas the $t$ in
$R_e(t)$ is $t^{\hgt}$ --- and $\hgt$ is precisely the statistic transposition
complements.  The instrument was reporting on a different referent.
\end{remark}

\section{The hyperbola as a fixed-point set}

On the parameter torus $(\C^\times)^2$ put
\[
  \varsigma(t,s)=(s,t),\qquad \iota(t,s)=(1/t,1/s),\qquad
  \sigma:=\varsigma\circ\iota:(t,s)\mapsto(1/s,1/t).
\]
These generate a group $\cong(\Z/2)^2$.  Note
$\operatorname{Fix}(\varsigma)=\{s=t\}$ and
\[
  \operatorname{Fix}(\sigma)=\{(t,s):s=1/t\}=\{ts=1\}.
\]

\begin{theorem}\label{thm:C}
For all $e,f\ge1$,
\[
  \omega\,C_{e,f}(t,s)\,\omega\;=\;t^{\,e-1}s^{\,f-1}\,C_{e,f}(1/t,1/s)
  \;=\;-\,t^{\,e-1}s^{\,f-1}\,C_{f,e}\bigl(\sigma(t,s)\bigr).
\]
In particular, for $e=f$ and $(t,s)$ on the hyperbola $ts=1$,
\[
  \boxed{\ \Ad(\omega)\,C_{e,e}(t,1/t)\;=\;-\,C_{e,e}(t,1/t).\ }
\]
\end{theorem}

\begin{proof}
The first equality is Theorem~\ref{thm:B} applied to each of the four products in
$C_{e,f}(t,s)=R_e(t)R_f(s)-R_f(s)R_e(t)$, using $\omega^2=\mathrm{id}$ to insert
$\omega\omega$ between the factors.  The second is
$C_{e,f}(a,b)=-C_{f,e}(b,a)$, valid by definition of the commutator, applied with
$(a,b)=(1/t,1/s)$ so that $(b,a)=(1/s,1/t)=\sigma(t,s)$.  For the boxed statement put
$e=f$ and $s=1/t$: then $\sigma(t,1/t)=(t,1/t)$ because $ts=1$ is the fixed locus of
$\sigma$, and $t^{e-1}s^{e-1}=(ts)^{e-1}=1$.
\end{proof}

So $ts=1$ \emph{is} the fixed-point set of a $\Z/2$ action --- of $\sigma$ on the
parameter torus, implemented on operators by $\Ad(\omega)$ together with the swap of
the two commutator slots.  Off the hyperbola, $\varsigma$ and $\iota$ are distinct
maps; on it they agree.  That is the sense, and the only sense, in which the locus is
the fixed-point set of an involution on these operators.

\section{Why the two-bead sector dies, and why the $\Z/2$ is not the reason}

Recall Theorem 4.1(2) of \texttt{2026-09-07-c2-Q96-order-over-sublattice.tex}: for
$|M\triangle M'|=4$, writing $M\setminus M'=\{b,c\}$, a legal assignment is an ordered
pair with $b+e\in M'\setminus M$, $c+f\in M'\setminus M$ and $b,c,b+e,c+f$ pairwise
distinct; with $P=\#(M\cap(b,b+e))$, $Q=\#(M\cap(c,c+f))$ and
\[
  k=\mathbf 1[\,b+e\in(c,c+f)\,]-\mathbf 1[\,b\in(c,c+f)\,],
\]
one has
$\langle M'|C_{e,f}(t,s)|M\rangle=\sum_{(b,c)\ \mathrm{legal}}t^{P-k}s^{Q}(1-(ts)^{k})$.

\begin{theorem}\label{thm:D}
For all $e,f\ge1$ the two-bead part of $C_{e,f}(t,s)$ vanishes identically on
$ts=1$.  Moreover, for $e=f$ the two-assignment matrix element factors as
\[
  \boxed{\ \langle M'|C_{e,e}(t,s)|M\rangle
  \;=\;\bigl(1-(ts)^{k}\bigr)\,\bigl(t^{P-k}s^{Q}-t^{Q}s^{P-k}\bigr).\ }
\]
\end{theorem}

\begin{proof}
Each summand carries the factor $1-(ts)^k$, which vanishes when $ts=1$ (for any
$k\in\Z$); the first claim follows termwise.  For the factorisation, $e=f$ gives
exactly two legal assignments, $(b,c)$ and $(c,b)$, with indices $k$ and $-k$ and
with $(P,Q)$ exchanged.  Writing $\alpha=P-k$, $\beta=Q$, the sum is
\[
  t^{\alpha}s^{\beta}\bigl(1-(ts)^k\bigr)+t^{\beta+k}s^{\alpha+k}\bigl(1-(ts)^{-k}\bigr)
  = \bigl(t^{\alpha}s^{\beta}-t^{\beta}s^{\alpha}\bigr)-(ts)^k
    \bigl(t^{\alpha}s^{\beta}-t^{\beta}s^{\alpha}\bigr),
\]
since $t^{\beta+k}s^{\alpha+k}(1-(ts)^{-k})=t^{\beta+k}s^{\alpha+k}-t^{\beta}s^{\alpha}$
and $t^{\alpha}s^{\beta}(ts)^k=t^{\alpha+k}s^{\beta+k}$.  Collecting gives the boxed
form.
\end{proof}

The factorisation makes the two exceptional loci visible as the two factors:
$1-(ts)^k$ is anti-invariant for $\sigma$ and cuts out $ts=1$; the alternant
$t^{\alpha}s^{\beta}-t^{\beta}s^{\alpha}$ is anti-invariant for $\varsigma$ and cuts
out $s=t$.  It also reproves the annotation recorded on 8 September: the element
vanishes identically iff $k=0$ or $\alpha=\beta$, i.e.\ iff $Q=P-k$.

\begin{proposition}[the mechanism, local on the Maya diagram]\label{prop:routes}
Let $b\to b+e$ and $c\to c+f$ be two bead moves on a Maya set $M$ with
$b,c,b+e,c+f$ pairwise distinct and $b+e,c+f\notin M$.  Let $h_e^0=\#(M\cap(b,b+e))$
and $h_f^0=\#(M\cap(c,c+f))$ be the heights in $M$, and let $h_e^{\mathrm{after}}$,
$h_f^{\mathrm{after}}$ be the heights of each move performed \emph{after} the other.
Then
\[
  h_e^{\mathrm{after}}-h_e^{0}=-k,\qquad h_f^{\mathrm{after}}-h_f^{0}=+k .
\]
Consequently the two orders carry weights
$t^{h_e^{0}}s^{h_f^{0}+k}$ (do $e$ first) and $t^{h_e^{0}-k}s^{h_f^{0}}$ (do $f$
first), whose ratio is exactly $(ts)^{k}$; the two orders are assigned \emph{equal}
weight precisely when $ts=1$, and the commutator then cancels in pairs.
\end{proposition}

\begin{remark}[a gloss of mine that was wrong, and the computation that caught it]
On first writing I asserted that reversing the order raises \emph{both} heights by
$k$.  That is false: the shifts are $-k$ and $+k$.  The conclusion survives ---
indeed the ratio is still $(ts)^k$ --- but for a different reason: each order
attributes the two heights to different variables, so the two $\pm k$ shifts appear
in the ratio with the \emph{same} sign.  A sweep of $30349$ disjoint move-pairs
returned $4343$ failures of the first statement and $0$ of the corrected one.
\end{remark}

\subsection*{The negative: the $\Z/2$ does not force the vanishing}

It is tempting to read Theorem~\ref{thm:C} as the cause of Theorem~\ref{thm:D}.  It
is not.  Membership of the $(-1)$-eigenspace of $\Ad(\omega)$ is not a vanishing
statement, and there is a live witness to the difference \emph{inside the same
operator}: the one-bead sector of $C_{e,f}$ satisfies the same eigenvalue equation on
$ts=1$ and is nonzero there.  Concretely, over $N\le4$:
\[
\begin{array}{lcc}
 (e,f) & \text{one-bead entries surviving on }ts=1 & \text{two-bead entries surviving}\\\hline
 (2,3) & 34/34 & 0/22\\
 (3,3) & 36/36 & 0/6\\
 (3,4) & 52/52 & 0/58
\end{array}
\]
The reason is visible in Proposition~\ref{prop:routes}: in the one-bead sector the
two routes move the \emph{same} bead ($b\to b+e\to b+e+f$ versus
$b\to b+f\to b+e+f$), so they are not a reordering of one pair of moves, and their
weight ratio is not of the form $(ts)^k$.  Order-reversal is available only when
there are two beads to reorder.

Hence the programme's stated hope --- that Q104 would make all four sightings of
$ts=1$ corollaries of one involution --- is refuted.  What survives is: (B) and (C)
explain why the locus is \emph{distinguished}; (D) explains the two-bead vanishing,
by a pairing on routes rather than a symmetry of operators.  The one-bead order-of-
vanishing results (Q105 Thms A and B) are \emph{not} consequences of the $\Z/2$ and
remain to be explained.

\section{Verification}

All computations are in \texttt{scratch/0909/}.  Two engines were used: an
abacus/Maya engine written this session, and the direct Young-diagram border-strip
engine \texttt{reviews/2026-09-08-selfreview-code/ribbon.py}, which is code-disjoint
from it and was previously cross-checked $764/764$ against Q105 Thm A.

\begin{enumerate}
\item \textbf{Theorem~\ref{thm:A}.} $1422$ pairs $(\mu,\text{strip})$ for
$e\in\{3,4,5\}$, $|\mu|\le9$: $0$ mismatches against $(r-1)e+\maj(R)$.  All $28$
ribbon shapes have non-constant $c$; all $28$ have $c-(r-1)e$ constant.
\item \textbf{Lemma~\ref{lem:transpose}.} $2670$ pairs, $e\le6$, $|\lambda|\le9$:
$0$ failures; every pair $(h,e-1-h)$ is realised.
\item \textbf{Theorem~\ref{thm:B}.} Verified as a matrix identity over $\Q(t)$ for
$e=1,\dots,5$ and $N=0,\dots,8$ (up to $101\times22$).  The two-parameter form of
Theorem~\ref{thm:C} verified for $(e,f)\in\{(2,2),(2,3),(3,3),(2,4),(3,4)\}$,
$N\le4$.
\item \textbf{Theorem~\ref{thm:C}, boxed form.} $\Ad(\omega)C_{e,e}=-C_{e,e}$ on
$ts=1$ verified for $e=2,3,4$, $N\le3$.
\item \textbf{Theorem~\ref{thm:D}.} Factorisation verified symbolically for
$P,Q\in\{0,\dots,5\}$, $k\in\{-2,\dots,2\}$.  Two-bead vanishing on $ts=1$ verified
for $(e,f)\in\{(1,2),(2,2),(2,3),(3,3),(3,4),(2,5),(4,4)\}$: $0$ survivors from
$131$ nonzero entries.
\item \textbf{Proposition~\ref{prop:routes}.} $30349$ randomly sampled disjoint
move-pairs, $e,f\le5$: $0$ failures.
\item \textbf{Controls.}  (a) $n(\mu+1^k)-n(\mu)$ is constant $=\binom k2$ for
$k=2,\dots,5$ over all $|\mu|\le10$ --- reproducing Garsia--Zabrocki's step ``using
2.10''.  (b) $\maj(1^k)=\binom k2$, so Theorem~\ref{thm:A} at $r=1$ recovers it.
(c) The same column shape added at $r\ge2$ gives $c\in\{6,9,12,\dots\}$: not
constant.  (d) The one-bead negative control \emph{moves} ($52/52$ against $0/58$),
so the two-bead statement is not degenerate.
\end{enumerate}

\section{Gaps and what is not claimed}

\begin{enumerate}
\item Theorem~\ref{thm:B} is proved for the operators $R_e(t)$ of
Convention~\ref{conv:main} only.  No claim is made about Jing's Hall--Littlewood
vertex operators $\Gamma^\pm(u)$, for which $\omega$-conjugation is $t$-blind
(Rozhkovskaya (4.1)--(4.2)); the two families are not identified here.
\item The scalar $t^{e-1}$ means $\Ad(\omega)$ implements $t\mapsto1/t$ only
projectively.  The normalisation $\widetilde R_e=t^{-(e-1)/2}R_e$ removes it at the
cost of leaving $\Q(t)$ for $\Q(t^{1/2})$.  Whether a $\Q(t)$-rational normalisation
exists is open and not addressed.
\item \textbf{The main gap.} Theorems~\ref{thm:C} and~\ref{thm:D} are both true and
are \emph{not} shown to be related.  I do not have a derivation of the two-bead
vanishing from the $\Z/2$, and the one-bead counterexample shows no such derivation
can be purely eigenvalue-theoretic.  A derivation would have to use the bead-number
grading, and I have not found one.
\item The Q105 order-of-vanishing statements at $t=-1$ (one-bead order exactly $1$;
vertex-side order $\equiv N \bmod 2$) are untouched by this paper.  They concern
$t=-1$, a point of $\operatorname{Fix}(\iota)$, not of $\operatorname{Fix}(\sigma)$;
whether the two loci interact is not investigated here.
\item Theorem~\ref{thm:A} assumes $\lambda/\mu$ connected.  The formula
\eqref{eq:split} --- $n(\lambda)-n(\mu)=\sum_{\text{cells}}(\mathrm{row}-1)$ --- holds
for any skew shape; only the identification of the second term with $\maj$ uses
connectedness.
\end{enumerate}

\end{document}
